\documentclass[12pt]{article}
\usepackage[T1]{fontenc}
\usepackage[utf8]{inputenc}
\usepackage[margin=1in]{geometry}
\usepackage{setspace}
\onehalfspacing
\usepackage[hidelinks]{hyperref}
\pagenumbering{gobble}

\begin{document}

\noindent Zhanbolat Kakishev\\
Nazarbayev University\\
Astana, Kazakhstan\\
\texttt{zhanbolat.kakishev@nu.edu.kz}

\vspace{1em}
\noindent 2 September 2026

\vspace{1em}
\noindent The Editors\\
\textit{Journal of International Economics}

\vspace{1.5em}
\noindent Dear Editors,

\vspace{0.8em}
\noindent I am pleased to submit ``Corridor, Not Factory: Trade Reorientation and the Missing
Investment Response in Kazakhstan, 2022--2025'' for consideration at the \textit{Journal of
International Economics}.

\vspace{0.8em}
\noindent After February 2022, part of the trade between Russia and its principal partners was
reallocated through third countries. The paper studies what one of those intermediaries,
Kazakhstan, retained from that reallocation, and whether it induced domestic investment. Using
product-level trade data (HS6, 2018--2025) I identify a 29-line ``surge basket'' in which
exports to Russia rise about tenfold and break sharply in May 2022; an externally compiled
product list, not selected on Kazakh outcomes, shows the same break. Propagating a
wholesale-and-freight margin through the input--output table implies that a rerouted dollar
generates on the order of a tenth of the domestic value added of a produced one, and I report
the full sensitivity of that figure to the chosen margin. Across three commercial deal
databases (Capital~IQ, PitchBook, Preqin) I find no investment response in the reoriented
product lines. I read the null through a three-gate decomposition---market access,
irreversibility, and capital-market institutions---and argue that a customs union with the
destination, which makes trade a close substitute for production, is close to sufficient for
the null on its own.

\vspace{0.8em}
\noindent The paper speaks to a question at the centre of the journal's scope: when a
trade-policy-driven demand shock is met by domestic production and when it merely passes
through. It contributes the first host-economy incidence analysis of the post-2022
reallocation---prior work, including \mbox{Chupilkin}, Javorcik and Plekhanov in the
\textit{European Economic Review}, establishes that the reallocation happened and its
scale---and it connects the entrep\^ot-margin tradition (Feenstra and Hanson on Hong Kong's
re-exports of Chinese goods) to the investment-under-uncertainty and institutional-voids
literatures through a single organising framework.

\vspace{0.8em}
\noindent The manuscript is original, is not under consideration elsewhere, and no part of it
has been published previously. It is single-authored. The trade analysis rests entirely on
public data (UN Comtrade, World Bank WDI, OECD ICIO, Kazakhstan BNS); the deal data are
licensed under academic terms and cannot be redistributed, but the full query specification
and the source deal identifiers are in the replication package, so the extract is
reconstructible by a licence holder. A complete replication package accompanies the paper. A
declaration of competing interest---the author is employed by the Astana International
Financial Centre---is on the title page and in the manuscript; it does not bear on the
empirical analysis.

\vspace{0.8em}
\noindent I suggest the following referees, none of whom has a conflict of interest to my
knowledge:
\begin{itemize}
\item \textbf{Julian Hinz} (Bielefeld University / Kiel Institute for the World Economy) ---
empirical work on sanctions and trade, including the measurement of trade re-routing through
third countries.
\item \textbf{Sonali Chowdhry} (DIW Berlin) --- export controls, trade diversion, and the
incidence of trade-policy shocks.
\item \textbf{R\'eka Juh\'asz} (University of British Columbia) --- when trade shocks induce a
domestic supply and industrial response, the paper's central question.
\end{itemize}
\noindent I have no objection to any potential referee.

\vspace{0.8em}
\noindent Thank you for considering the paper.

\vspace{1.5em}
\noindent Sincerely,\\[1.5em]
Zhanbolat Kakishev

\end{document}
